% Where to look: LLM-free, CPU-only code localisation for coding agents, measured on holdouts.
%
% Source of every number: docs/research/contributions-log.md (sections cited in comments as "log sec. x").
% Competitor numbers: docs/research/related-work-2026.md (PR #146).
% Build: `make` (one column) or `make twocol`. Two-column mode is selected by defining
% \twocolumnmode to 1 before this file is read (see the Makefile).
\providecommand{\twocolumnmode}{0}
\if1\twocolumnmode
  \documentclass[10pt,twocolumn]{article}
  \PassOptionsToPackage{margin=0.75in}{geometry}
\else
  \documentclass[11pt]{article}
  \PassOptionsToPackage{margin=1in}{geometry}
\fi

\usepackage[T1]{fontenc}
\usepackage[utf8]{inputenc}
\usepackage{lmodern}
\usepackage{geometry}
\usepackage{microtype}
\usepackage{amsmath,amssymb}
\usepackage{booktabs}
\usepackage{graphicx}
\usepackage{xcolor}
\usepackage{url}
\usepackage[numbers,sort&compress]{natbib}
\usepackage[colorlinks=true,linkcolor=blue!50!black,citecolor=blue!50!black,urlcolor=blue!50!black]{hyperref}

% Unmeasured or unconfirmed content. `make check` lists every occurrence; none may remain at submission.
\newcommand{\todo}[1]{\textcolor{red}{\textbf{[TODO:~#1]}}}
\newcommand{\nm}{NeuroMesh}
\newcommand{\accat}[1]{Acc@#1}

\title{Where to Look: LLM-Free, CPU-Only Code Localisation\\for Coding Agents, Measured on Holdouts}
\author{Parsa Karkooti\thanks{\todo{affiliation and co-authors}}}
\date{}

\begin{document}
\maketitle

% ---------------------------------------------------------------------------------------------
\begin{abstract}
Coding agents spend much of their budget finding the code a task needs. Strong localisers today
run an LLM agent over the repository, or embed every function of it with a trained retriever,
usually on a GPU. We ask how far a local engine gets with neither: a per-project code graph,
field-weighted lexical ranking, a definition-level ranking of long reports, report hygiene, and the
report's own traceback, all on a laptop CPU in a few seconds per issue. Design decisions were
tuned on SWE-bench dev only (the last one also on a SWE-Gym split, after the released engine had
been run on it once), and each released version was run once on held-out data. On 276 held-out SWE-bench Lite issues the engine reaches file-level
\accat{1}/3/5 of 0.507/0.717/0.750 (BM25 0.301/0.507/0.587) at 2.9\,s per issue including a cold
index; on the 407 SWE-bench Verified issues outside Lite 0.435/0.671/0.732 (BM25
0.194/0.373/0.477); on 280 SWE-Gym issues from seven repositories it never saw, file \accat{1}
0.279/0.271 on the two halves (BM25 0.121/0.121). At function level it lists every edited function
among its first five on 0.460 of LocAgent's 274 Lite instances, and 0.485 with an opt-in vote that
embeds only the engine's top 30 definitions on CPU (paired $+0.026$, 95\% CI [0.004, 0.051]); on
Verified the same vote lifts function \accat{5} from 0.429 to 0.458 ($+0.028$ [0.009, 0.048]). This
is below code-trained dense retrievers --- SweRankEmbed-Small reaches 0.664 file \accat{1} and
0.631 function \accat{5} on the same instances --- and well below LLM agents and rerankers. Our
contribution is the other corner of the trade-off: no training, no GPU, no whole-repository
embedding index, zero tokens, a holdout protocol that caught two results looked-at data had
suggested, and a dozen negative results with numbers.
\end{abstract}

% ---------------------------------------------------------------------------------------------
\section{Introduction}
\label{sec:intro}

A coding agent asked to fix an issue first has to find where the fix goes. On SWE-bench
\citep{jimenez2024swebench} this step, \emph{localisation}, is where agent pipelines spend much of
their budget: Agentless \citep{xia2025agentless} prompts a model with the repository tree, then with
file skeletons, then with code; LocAgent \citep{chen2025locagent}, SWE-agent \citep{yang2024sweagent}
and OpenHands \citep{wang2025openhands} run multi-step LLM searches. The strongest model-free
alternatives embed every function of the repository with a trained encoder
\citep{suresh2025cornstack,reddy2026swerank}. A 137M-parameter encoder is small for a GPU, but on a
laptop CPU we measured 2.6 chunks per second for a comparable model (\S\ref{sec:negative}): hours per
large repository, repeated whenever the code changes. Local-first assistants such as Aider's RepoMap
\citep{aider_repomap} ship structural or lexical retrieval but are not measured on localisation
benchmarks.

We ask a narrow question: \emph{how far does a local engine get with no LLM and no GPU, if every
design decision is measured on a tuning split and confirmed once on held-out data?} The engine,
\nm{} \citep{neuromesh}, builds a per-project code graph in seconds, ranks files with field-weighted
BM25 \citep{robertson2004bm25f}, ranks \emph{definitions} for long reports, removes issue-template
scaffolding from the query, and reads the report's own pointers --- the files and functions its
traceback runs through. It returns a ranked list of files and of functions (\texttt{where\_to\_look},
\texttt{functions\_to\_look}) that an agent opens first.

The answer is \emph{not} state of the art, and we say so up front. On LocAgent's 274 Lite instances
the engine is 0.03 below CodeRankEmbed and 0.16 below SweRankEmbed-Small at file \accat{1}, and
0.27 below LocAgent with Claude-3.5 at function \accat{5} (Table~\ref{tab:lite274}). What it offers
is a measured point no published method occupies: no learned model at query time, no
whole-repository index beyond a code graph, zero tokens, 2.9\,s per issue on a laptop, and numbers
from holdouts run once.

\paragraph{Contributions.}
\begin{enumerate}
  \item \textbf{A holdout protocol for retrieval engines} (\S\ref{sec:protocol}): gold locked by
  commit before any run, two tuning splits disjoint in repositories from every evaluation set, a
  pre-set acceptance bar, one run per released version on each holdout, and per-set labels for what
  has been looked at. It caught two results that looked-at sets suggested (a cross-encoder reranker
  and a neighbour vote) and confirmed one on a fresh repository.
  \item \textbf{An LLM-free localiser} reaching file \accat{1} 0.507 on 276 held-out Lite issues (BM25
  0.301), 0.435 on 407 Verified issues outside Lite (BM25 0.194) and $2.2$--$2.3\times$ BM25 on seven unseen
  SWE-Gym repositories, at 2.9\,s per issue on a laptop CPU (\S\ref{sec:results}).
  \item \textbf{A cheap dense short-list vote}: embedding only the engine's top 30 definitions lifts
  function \accat{5} on Lite from 0.460 to 0.485 and on Verified from 0.429 to 0.458, with the
  decision committed before both holdout runs
  (\S\ref{sec:dense}).
  \item \textbf{Ablations, costs and negative results with numbers} (\S\ref{sec:ablation},
  \S\ref{sec:negative}): a dozen ideas from the bug-localisation and retrieval literature --- commit
  history, query expansion, cross-encoder reranking, graph expansion, a learned reranker, a small
  local LLM judge --- that did not survive the protocol, and why.
\end{enumerate}

% ---------------------------------------------------------------------------------------------
\section{Related work}
\label{sec:related}

\paragraph{LLM localisers.}
Agentless \citep{xia2025agentless} narrows hierarchically (repository tree, files, class and
function skeletons, lines) with one model call per level plus sampling and voting; with GPT-4o it
covers the gold file in 69.7\% and the gold function in 52.0\% of Lite issues at \$0.70 per issue.
SWE-agent \citep{yang2024sweagent}, OpenHands \citep{wang2025openhands} and MoatlessTools
\citep{moatless} search through shell or retrieval tools inside an agent loop. LocAgent
\citep{chen2025locagent} indexes a heterogeneous code graph (files, classes, functions;
contain/import/invoke/inherit edges) that an LLM agent traverses with entity-level BM25. It also
fixed the reference protocol we use: 274 Lite instances whose patch edits an existing function, and
\accat{k} requiring all gold locations in the top $k$. With Claude-3.5 it reaches file \accat{1}/5
0.777/0.942 and function \accat{5} 0.734 at \$0.66 per issue. Later systems refine the search
policy: iterative graph search (CoSIL, \citealp{jiang2025cosil}), priority scheduling and context
pruning (OrcaLoca, \citealp{yu2025orcaloca}), repository knowledge graphs (KGCompass,
\citealp{yang2025kgcompass}), repository memory (RepoMem, \citealp{wang2025repomem}) and code graphs as
a plug-in (RepoGraph, \citealp{ouyang2025repograph}). Another line trains small open models for
localisation with reinforcement learning \citep{ma2026tooltrain,sutawika2026codescout,xu2026fusesearch}.
With frontier models the numbers keep rising: SWE-Adept \citep{he2026sweadept} reports function
\accat{5} 0.878 on the 274 instances with Claude-Sonnet-4.5 at 348k tokens per issue, and in the same
paper LocAgent re-run with GPT-5.2 scores 0.646 against the 0.734 originally reported --- agent
numbers move by about 0.1 with the model and the run. SHERLOC \citep{tamoyan2026sherloc} adds
structured diagnosis with a reasoning model. Our engine borrows LocAgent's entity-level content
ranking and Agentless's file-to-function narrowing, without the model.

\paragraph{Dense retrieval and learned rerankers.}
CodeRankEmbed (137M; \citealp{suresh2025cornstack}) ranks function chunks by cosine with the issue;
on the 274 instances it reaches file \accat{1} 0.526 and function \accat{5} 0.518, and a 7B
listwise reranker lifts the latter to 0.675. SweRank \citep{reddy2026swerank} fine-tunes retrievers
and rerankers on issue-localisation data mined from GitHub: SweRankEmbed-Small (137M) reaches
0.664/0.631 (file \accat{1}/function \accat{5}) with no LLM at query time, and
SweRankEmbed-Large with a 32B reranker 0.832/0.814 at \$0.015 per issue on a GPU; SweRank+
\citep{reddy2025swerankplus} extends this to several languages (function \accat{5} 0.854). CoRet
\citep{fehr2025coret} adds call-graph context to a 130M retriever, SWE-Fixer \citep{xie2025swefixer}
selects files from a BM25 short list with a fine-tuned 7B model, and BLAZE
\citep{chakraborty2025blaze} fine-tunes CodeSage \citep{zhang2024codesage} for cross-project bug
localisation. These are the strongest non-agent localisers. All require training data close to the
task and an embedding of every function in the repository; none reports CPU indexing cost. We
measure that cost, and a variant that embeds only the engine's short list (\S\ref{sec:dense}).

\paragraph{Classic bug localisation.}
IR-based bug localisation uses report text, history and structure: BugLocator
\citep{zhou2012buglocator} combines a revised VSM with similar fixed bugs, Locus
\citep{wen2016locus} ranks change hunks, and \citet{rahman2018query} reformulate the query. We test
the history and reformulation (RM3, \citealp{abduljaleel2004rm3}) ideas under our protocol; both are
negative (\S\ref{sec:negative}). MULocBench \citep{zhang2025mulocbench} shows that issues outside
code (configuration, documentation, build) are much harder --- LocAgent with Claude-3.5 reaches file
\accat{5} 0.352 there --- which motivates our plain-language holdouts (\S\ref{sec:plain}).

\paragraph{Ranking machinery.}
We use BM25F \citep{robertson2004bm25f} and reciprocal rank fusion (RRF;
\citealp{cormack2009rrf}); \citet{bruch2023fusion} analyse convex lexical/dense fusion.

\paragraph{Positioning.}
No published method in this space is free of learned models at query time, and none reports a
pre-registered holdout. Table~\ref{tab:lite274} places our engine among the methods measured on the
same 274 instances with the same metric.

% ---------------------------------------------------------------------------------------------
\section{System}
\label{sec:system}

\paragraph{Index.}
A per-project graph of files, definitions (functions, methods, classes), calls and imports, built by
tree-sitter parsers and cached on disk. Link resolution was the cold-start bottleneck: on a django
checkout (3.5k files) four fixes to repeated path normalisation and per-relation neighbour walks
brought the first packet from 3\,min\,20\,s to 26\,s without changing the graph (log~\S8.2).

\paragraph{Questions.}
A prompt's anchors (identifiers, paths) seed spreading activation over the graph; the activated
nodes form a token-bounded \emph{packet}. Files are ranked by BM25F over four fields: path, symbol
names, comments and body.

\paragraph{Reports.}
For prompts of at least 60 words the engine treats the prompt as an issue report. (i)~\emph{Hygiene}
removes issue-template scaffolding (``search before asking'', checklists, links, headings) and keeps
code-like tokens that word tokenisers drop (rule codes such as \texttt{L031}). (ii)~A
\emph{definition-level} BM25 ranks every definition's text, with the report title counted three
times; files inherit the rank of their best definition. (iii)~\emph{Named files}: traceback frames
(deepest first), then paths and unique file names in prose, resolved against the index by longest
unique path suffix. The packet order, the definition ranking and the named files are fused by RRF
(named files at weight~2) into the file list \texttt{where\_to\_look}.

\paragraph{Short questions.}
For short prompts the file list is the RRF of the packet order and a whole-question ranking (S1);
the packet's activation order alone puts hub files first.

\paragraph{Functions to look at.}
For every prompt the definition ranking (classes removed) is fused by RRF with each definition's file
rank --- Agentless's file-to-function step without an LLM --- and with the functions a traceback runs
through: each frame is resolved to an indexed file, then to the definition named in the frame (or
the innermost one containing its line), test frames skipped, voting at double weight deepest-first.

\paragraph{Opt-in dense short-list vote.}
With \texttt{NM\_DENSE=1}, jina-embeddings-v2-base-code \citep{gunther2023jina} (int8 ONNX, CPU)
embeds the issue and the engine's top 30 definitions (text from the base commit, cut at 512 tokens);
their cosine order joins the function-list fusion by RRF ($k=60$, weight $0.5$). The file list is not
changed. It costs about 6\,s per issue on CPU and 31 embeddings (30 definitions and the issue), instead of an embedding of every
function in the repository.

% ---------------------------------------------------------------------------------------------
\section{Evaluation protocol}
\label{sec:protocol}

\paragraph{Tuning splits.}
\emph{SWE-bench dev} (225 issues from pvlib, pydicom, sqlfluff, astroid, pyvista and marshmallow ---
no repository shared with Lite or Verified), in halves dev-fast (59, ten per repository) and dev-rest
(166). Because one dev-fast issue is worth 0.017, dev cannot resolve differences of 0.02, and its
function gold is unlike Lite's; we therefore added \emph{SWE-Gym} \citep{pan2025swegym}: 280 issues
from seven repositories disjoint from Lite, Verified and dev (pydantic, hydra, conan, dvc, dask,
moto, mypy), at most 40 per repository, patches that edit only existing non-test Python files,
sampled with a fixed seed by a script committed before any run, in halves gym-a and gym-b.

\paragraph{Acceptance bar.}
A change is kept only if file \accat{1} or function \accat{5} rises by at least 0.02 in \emph{both}
halves of a tuning split, a bar fixed before each measurement. Every idea is first a Python
prototype over stored engine outputs, then an engine port re-measured on the full split.

\paragraph{Holdouts.}
SWE-bench Lite test (300) and Verified (500), each run once per released version with the release
binary and never inspected per instance. The 24 Lite instances looked at in an early session are
excluded from the strict row (276). LocAgent's 274-instance subset (patches that edit an existing
function) was rebuilt from its definition and matches exactly; it gives the head-to-head. 407
Verified issues are outside Lite. The released v1.4.0 had never seen the SWE-Gym repositories, so its
single run there is a third holdout, reported before any tuning on SWE-Gym.

\paragraph{Metrics.}
File-level \accat{k}: all gold files in the top $k$. Function-level \accat{k}: all edited functions in
the top $k$, where function gold is the innermost function containing a removed line or insertion
point of the reference patch at the base commit; an instance without a function list counts as a
miss. 95\% bootstrap confidence intervals \citep{efron1993bootstrap}; paired comparisons report the
bootstrap CI of the difference, issues gained and lost, and a two-sided sign test.

\paragraph{Plain-language questions.}
Five holdouts (ripgrep/Rust, click/Python, cobra/Go, axios/JavaScript, jsoup/Java), 12 questions
each, gold written from source and locked by commit before the first run; metric R@3 of the file
list.

% ---------------------------------------------------------------------------------------------
\section{Results}
\label{sec:results}

% Head-to-head on LocAgent's 274-instance SWE-bench Lite subset (same instances, same metric).
% Ours: contributions-log sec. 8.7, sec. 8.13, sec. 8.17, sec. 8.18. Others: related-work-2026.md sec. 1
% (LocAgent Table 4, SweRank Table 1/5), rounded to three decimals.
\begin{table*}[t]
\centering
\small
\caption{SWE-bench Lite, LocAgent's 274 instances (patch edits an existing function); Acc@$k$ = all
gold locations in the top $k$. Rows marked \dag{} are copied from LocAgent~\citep[Table~4]{chen2025locagent}
and SweRank~\citep[Tables~1 and~5]{reddy2026swerank}. ``Model'' is what runs at query time; ``index'' is
what must be computed per repository before the first query. Our rows are a single run of a released
binary on a holdout (\S\ref{sec:protocol}); 95\% bootstrap CIs in brackets. The dense vote changes only
the function list, so its file columns equal the engine's.}
\label{tab:lite274}
\setlength{\tabcolsep}{4pt}
\resizebox{\textwidth}{!}{%
\begin{tabular}{@{}llcccccl@{}}
\toprule
 & & \multicolumn{3}{c}{file Acc@$k$} & \multicolumn{2}{c}{function Acc@$k$} & \\
\cmidrule(lr){3-5}\cmidrule(lr){6-7}
method & model at query time / index & @1 & @3 & @5 & @5 & @10 & cost per issue \\
\midrule
\multicolumn{8}{@{}l}{\emph{No learned model}}\\
BM25\dag & -- / inverted index & 0.387 & 0.518 & 0.617 & 0.318 & 0.369 & -- \\
BM25 (our implementation) & -- / inverted index & 0.299 & 0.522 & 0.606 & -- & -- & -- \\
\textbf{NeuroMesh v1.4.0 (ours)} & -- / code graph, 26\,s cold (django) & \textbf{0.500} & \textbf{0.723} & \textbf{0.755} & \textbf{0.460} & \textbf{0.540} & CPU, p50 2.9\,s, 0 tokens \\
 & & {\scriptsize[0.44,0.56]} & & & {\scriptsize[0.40,0.52]} & & \\
\midrule
\multicolumn{8}{@{}l}{\emph{Short-list embedding (30 definitions per issue, no repository embedding)}}\\
\textbf{NeuroMesh v1.4.0 + dense vote (ours)} & 161M encoder, CPU / code graph & 0.500 & 0.723 & 0.755 & \textbf{0.485} & \textbf{0.555} & CPU, +$\approx$6\,s p50 \\
 & & & & & {\scriptsize[0.43,0.54]} & & \\
\midrule
\multicolumn{8}{@{}l}{\emph{Dense retrievers (embed every function of the repository)}}\\
CodeSage-large-v2 (1.3B)~\citep{zhang2024codesage}\dag & -- / all functions & 0.478 & 0.693 & 0.781 & 0.339 & 0.445 & GPU in practice \\
Jina-Code-v2 (161M)~\citep{gunther2023jina}\dag & -- / all functions & 0.434 & 0.712 & 0.803 & 0.423 & 0.522 & -- \\
E5-base-v2 (110M)~\citep{wang2022e5}\dag & -- / all functions & 0.496 & 0.745 & 0.803 & 0.394 & 0.511 & -- \\
CodeRankEmbed (137M)\dag & -- / all functions & 0.526 & 0.777 & 0.847 & 0.518 & 0.588 & -- \\
SweRankEmbed-Small (137M)\dag & -- / all functions & 0.664 & 0.865 & 0.909 & 0.631 & 0.745 & -- \\
SweRankEmbed-Large (7B)\dag & -- / all functions & 0.726 & 0.912 & 0.942 & 0.719 & 0.821 & GPU \\
\midrule
\multicolumn{8}{@{}l}{\emph{Retriever + LLM reranker}}\\
SweRankEmbed-L + SweRankLLM-L\dag & 32B reranker, GPU / all functions & 0.832 & 0.949 & 0.960 & 0.814 & 0.887 & \$0.015 \\
\midrule
\multicolumn{8}{@{}l}{\emph{LLM agents}}\\
SWE-agent + GPT-4o\dag & API LLM / -- & 0.573 & 0.650 & 0.690 & 0.460 & 0.464 & \$0.56 \\
OpenHands + GPT-4o\dag & API LLM / -- & 0.610 & 0.719 & 0.737 & 0.496 & 0.504 & \$0.83 \\
Agentless + GPT-4o\dag & API LLM / -- & 0.672 & 0.745 & 0.745 & 0.555 & 0.555 & -- \\
Agentless + Claude-3.5\dag & API LLM / -- & 0.726 & 0.792 & 0.796 & 0.588 & 0.588 & -- \\
LocAgent + Qwen2.5-7B (ft)\dag & 7B LLM, GPU / code graph & 0.708 & 0.847 & 0.883 & 0.642 & 0.715 & \$0.05 \\
LocAgent + Claude-3.5\dag & API LLM / code graph & 0.777 & 0.920 & 0.942 & 0.734 & 0.774 & \$0.66 \\
\bottomrule
\end{tabular}}
\end{table*}

% File and function level on every SWE-bench-style set, BM25 vs engine.
% contributions-log sec. 8.13 (file, v1.3.0 = v1.4.0 at file level per sec. 8.17), sec. 8.17 (function, v1.4.0),
% sec. 8.18 (SWE-Gym), paper-draft sec. 5.2 (dev).
\begin{table*}[t]
\centering
\small
\caption{File- and function-level localisation on four holdouts and the tuning split. Engine =
NeuroMesh v1.4.0, one run per set (file lists are identical to v1.3.0 on every instance). Function
columns count only instances with function gold ($n_f$); an instance with no function list counts as a
miss. 95\% bootstrap CI for file Acc@1 in brackets.}
\label{tab:holdouts}
\setlength{\tabcolsep}{4pt}
\resizebox{\textwidth}{!}{%
\begin{tabular}{@{}lr cccc cccc r ccc@{}}
\toprule
 & & \multicolumn{4}{c}{BM25, file Acc@$k$} & \multicolumn{4}{c}{engine, file Acc@$k$} & & \multicolumn{3}{c}{engine, function Acc@$k$}\\
\cmidrule(lr){3-6}\cmidrule(lr){7-10}\cmidrule(l){12-14}
set & $n$ & @1 & @3 & @5 & @10 & @1 & @3 & @5 & @10 & $n_f$ & @1 & @5 & @10\\
\midrule
\multicolumn{14}{@{}l}{\emph{Holdouts (never tuned on; run once per released version)}}\\
Lite, strict & 276 & 0.301 & 0.507 & 0.587 & 0.721 & \textbf{0.507} {\scriptsize[0.45,0.57]} & 0.717 & 0.750 & 0.815 & 250 & 0.228 & 0.472 & 0.556\\
Lite, LocAgent subset & 274 & 0.299 & 0.522 & 0.606 & 0.734 & \textbf{0.500} {\scriptsize[0.44,0.56]} & 0.723 & 0.755 & 0.828 & 274 & 0.223 & 0.460 & 0.540\\
Verified, all & 500 & 0.216 & 0.392 & 0.490 & 0.642 & \textbf{0.446} {\scriptsize[0.40,0.49]} & 0.680 & 0.736 & 0.814 & 459 & 0.198 & 0.429 & 0.497\\
Verified, not in Lite & 407 & 0.194 & 0.373 & 0.477 & 0.622 & \textbf{0.435} {\scriptsize[0.38,0.48]} & 0.671 & 0.732 & 0.811 & 375 & 0.189 & 0.397 & 0.464\\
SWE-Gym, gym-a & 140 & 0.121 & 0.257 & 0.329 & 0.407 & \textbf{0.279} {\scriptsize[0.21,0.36]} & 0.450 & 0.536 & 0.650 & -- & 0.152 & 0.258 & 0.333\\
SWE-Gym, gym-b & 140 & 0.121 & 0.271 & 0.307 & 0.436 & \textbf{0.271} {\scriptsize[0.20,0.35]} & 0.414 & 0.493 & 0.600 & -- & 0.112 & 0.254 & 0.366\\
\midrule
\multicolumn{14}{@{}l}{\emph{Tuning split (looked at)}}\\
SWE-bench dev & 225 & 0.160 & 0.347 & 0.427 & 0.538 & 0.308 & 0.527 & 0.621 & 0.692 & 210 & 0.190 & 0.324 & 0.390\\
\bottomrule
\end{tabular}}
\par\smallskip
{\footnotesize SWE-Gym: 266 of the 280 issues have function gold; the per-half split of $n_f$ is not
recorded in the log.}
\end{table*}


\subsection{Head-to-head on LocAgent's 274 Lite instances}
\label{sec:h2h}

Table~\ref{tab:lite274} compares the engine with every method reported on the same instances and
metric. Without any learned model it is $+0.11$ above LocAgent's BM25 at file \accat{1} (0.500 vs
0.387) and $+0.14$ at function \accat{5} (0.460 vs 0.318). It is above Jina-Code-v2 at file
\accat{1} (0.434), level with E5-base (0.496), and below CodeRankEmbed (0.526) and SweRankEmbed-Small
(0.664). At file \accat{5} it is below every dense retriever (0.755 vs 0.781--0.909). At function
\accat{5} it is above E5-base, Jina-Code-v2 and CodeSage-large (0.394, 0.423, 0.339), equal to
SWE-agent with GPT-4o (0.460), and below CodeRankEmbed (0.518), SweRankEmbed-Small (0.631) and every
LLM agent on Claude-3.5. Function \accat{1} is 0.223 (not reported by the other papers).

The 24 instances looked at during development are all among the 274; without them (250 instances)
the v1.2.0 engine scored 0.488/0.716/0.756 and BM25 0.304/0.512/0.592, so the looked-at instances do
not flatter the result (log~\S8.7).

\subsection{Three holdouts}
\label{sec:holdouts}

Table~\ref{tab:holdouts} gives every SWE-bench-style set. The gain over BM25 holds across all of
them: file \accat{1} $+0.21$ on Lite strict, $+0.24$ on Verified outside Lite, and $2.2$--$2.3\times$ on
the two SWE-Gym halves --- seven repositories the engine was never developed on. On dev the gain was
$+0.15$; the holdout gains are larger, not smaller, which is the expected sign if the tuning did not
overfit dev.

Function-level gains from dev transfer \emph{larger} than predicted. The traceback vote and function
lists for short reports (v1.3.0 $\to$ v1.4.0) raised dev function \accat{5} by $+0.034$
(0.290 $\to$ 0.324) but Lite-274 by $+0.066$ (0.394 $\to$ 0.460) and Verified by $+0.083$ (0.346
$\to$ 0.429). The inputs do not explain this --- tracebacks appear in 20\%, 20\% and 14\% of dev,
Lite and Verified issues, reports under 60 words in 10\%, 12\% and 14\% --- but the gold does: one
edited function in 45\% of dev issues (mean 3.7) against 86\% of Lite (mean 1.17) and 72\% of
Verified (mean 1.87). When every edited function must be in the top $k$, a correctly placed function
pays off far more often on single-function gold. SWE-Gym's gold (mean 3.4 functions, one function in
45\%) resembles dev, which is why its function numbers are low.

\subsection{Dense short-list vote}
\label{sec:dense}

% Dense short-list vote: dev near miss (sec. 8.16), SWE-Gym pass + pooled test, decision locked in
% 608acb4, then Lite once (sec. 8.18). Verified still being measured.
\begin{table}[t]
\centering
\small
\caption{Opt-in dense short-list vote (jina-code-v2, int8, CPU; the issue and the engine's top-30
definitions are embedded, their cosine order votes by RRF at $w=0.5$). Function Acc@$k$, engine $\to$
engine + vote. The weight and $k$ were committed before the Lite and Verified runs. Paired
differences: bootstrap 95\% CI, issues gained/lost at @5, two-sided sign test.}
\label{tab:dense}
\setlength{\tabcolsep}{3pt}
\begin{tabular}{@{}lrccc@{}}
\toprule
set & $n$ & @1 & @5 & @10\\
\midrule
dev-fast (tuning) & 52 & 0.173\,$\to$\,0.212 & 0.308\,$\to$\,0.327 & --\\
dev-rest (tuning) & 158 & 0.196\,$\to$\,0.215 & 0.329\,$\to$\,0.342 & --\\
gym-a (tuning) & 140 & 0.152\,$\to$\,0.167 & 0.258\,$\to$\,0.288 & 0.333\,$\to$\,0.364\\
gym-b (tuning) & 140 & 0.112\,$\to$\,0.134 & 0.254\,$\to$\,0.313 & 0.366\,$\to$\,0.388\\
\addlinespace
\multicolumn{5}{@{}l}{dev + gym pooled, $n=476$: @5 0.286\,$\to$\,0.317, $+0.032$ [0.015,\,0.050],}\\
\multicolumn{5}{@{}l}{\quad 17 gained / 2 lost, $p=0.0007$}\\
\midrule
\multicolumn{5}{@{}l}{\emph{Holdout, run once with the locked decision}}\\
Lite, LocAgent subset & 274 & 0.223\,$\to$\,\textbf{0.252} & 0.460\,$\to$\,\textbf{0.485} & 0.540\,$\to$\,\textbf{0.555}\\
Lite, strict & 250 & 0.228\,$\to$\,0.252 & 0.472\,$\to$\,0.496 & 0.556\,$\to$\,0.568\\
\addlinespace
\multicolumn{5}{@{}l}{Lite 274 paired: @5 $+0.026$ [0.004,\,0.051], 9 gained / 2 lost, $p=0.065$;}\\
\multicolumn{5}{@{}l}{\quad @1 $+0.029$ (18 / 10)}\\
\addlinespace
Verified & 459 & \multicolumn{3}{c}{\todo{locked-decision run, being measured}}\\
\bottomrule
\end{tabular}
\par\smallskip
{\footnotesize Lite numbers come from the Python prototype applied to the stored v1.4.0 run. The Rust
port (\texttt{NM\_DENSE=1}) reproduces the prototype order up to near-tie swaps; on SWE-Gym it gives
@5 0.295 (gym-a) and 0.299 (gym-b).}
\end{table}


The dense vote of \S\ref{sec:system} first came up as a near miss on dev: positive in almost every
cell but below the bar on dev-fast (file \accat{1} $+0.017$, one issue; function \accat{5} $+0.019$
and $+0.013$). Re-measured unchanged on SWE-Gym, it passed (function \accat{5} $+0.030$ and $+0.059$;
Table~\ref{tab:dense}). Pooled over dev and SWE-Gym (476 issues) function \accat{5} rose from 0.286
to 0.317, $+0.032$ [0.015, 0.050], 17 issues gained and 2 lost. The weight and short-list size were
committed before the Lite run. On Lite-274, run once, function \accat{1}/5/10 went from
0.223/0.460/0.540 to 0.252/0.485/0.555: $+0.026$ at \accat{5}, CI [0.004, 0.051], 9 gained and 2
lost. The CI excludes zero; with only 11 discordant issues the sign test does not ($p=0.065$). The
effect is the same direction and size as on the tuning splits. It narrows the gap to CodeRankEmbed
(0.518) from $-0.058$ to $-0.033$ and leaves a gap of $-0.146$ to SweRankEmbed-Small, with 30
embeddings per issue on CPU instead of every function of the repository. On Verified (459 issues
with function gold), also run once with the locked decision, function \accat{1}/5/10 went from
0.198/0.429/0.497 to 0.237/0.458/0.514: $+0.028$ at \accat{5}, CI [0.009, 0.048], 17 gained and 4
lost ($p=0.007$), and $+0.039$ at \accat{1} (28/10, $p=0.005$); on the 375 outside Lite \accat{5}
rose from 0.397 to 0.435. Every split --- dev, both SWE-Gym halves, Lite and Verified --- moves the
same way by $+0.02$ to $+0.06$ at function \accat{5}, and the effect predicted on tuning data
($+0.032$) is what both holdouts show.

The vote also re-ranks files: the dense order alone puts the gold file first more often than the
engine on dev (0.356 vs 0.307 file \accat{1}) while losing depth, but the file vote did not clear the
bar on gym-b ($+0.000$), so the file list is unchanged.

\subsection{Ablation}
\label{sec:ablation}

% Ablation, v1.2.0 engine, Lite strict holdout (276); contributions-log sec. 8.11.
\begin{table}[t]
\centering
\small
\caption{Ablation on the Lite strict holdout (276 issues), one run per removed component, engine
v1.2.0 (before the named-files and traceback votes). File Acc@$k$.}
\label{tab:ablation}
\begin{tabular}{@{}lcccc@{}}
\toprule
removed & @1 & @3 & @5 & @10\\
\midrule
nothing (v1.2.0) & 0.486 & 0.710 & 0.750 & 0.804\\
report hygiene + code tokens & 0.467 & 0.707 & 0.743 & 0.801\\
definition-level ranking & \textbf{0.388} & \textbf{0.569} & \textbf{0.652} & 0.764\\
localisation list (packet order only) & 0.366 & 0.536 & 0.558 & --\\
\midrule
plain BM25, for reference & 0.301 & 0.507 & 0.587 & 0.721\\
\bottomrule
\end{tabular}
\end{table}


Table~\ref{tab:ablation} removes one component at a time on the Lite strict holdout (engine v1.2.0,
before the named-file and traceback votes). The definition-level ranking --- LocAgent's
entity-content BM25, made LLM-free --- carries most of the gain over the packet ($+0.10$ \accat{1},
$+0.14$ \accat{3}). Report hygiene is small on Lite ($-0.019$ \accat{1} when removed); its dev gains
came mostly from sqlfluff's template-heavy issues. The localisation list itself matters most: the
packet order alone scores 0.366 \accat{1}, because a packet holds one to five files and caps
\accat{5}. The named-files vote was added after this ablation; its effect is the v1.2.0 $\to$ v1.3.0
step on Verified outside Lite, file \accat{1} 0.392 $\to$ 0.437 on the same 403 instances, against
$+0.049$ predicted on dev (log~\S8.13).

\subsection{Plain-language questions}
\label{sec:plain}

% Plain-language holdouts, R@3 of the localisation list; contributions-log sec. 8.6, sec. 8.16.
\begin{table}[t]
\centering
\small
\caption{Plain-language questions (12 per repository, gold written from the source and committed
before the first run), R@3 of the localisation list. S1 = fusing the packet order with the
whole-question ranking for short questions; it was chosen on the first three sets.}
\label{tab:plain}
\begin{tabular}{@{}lccc@{}}
\toprule
set (language) & BM25 & before S1 & S1\\
\midrule
ripgrep (Rust), looked at & 0.500 & 0.208 & 0.500\\
click (Python), looked at & 0.750 & 0.750 & 0.917\\
cobra (Go), looked at & 0.875 & 0.792 & 0.875\\
axios (JS), fresh at its run & 0.583 & 0.833 & 0.833\\
\textbf{jsoup (Java), fresh, run once} & 0.667 & 0.417 & \textbf{0.750}\\
\bottomrule
\end{tabular}
\end{table}


Issue reports are long and name code. Developers' questions to an assistant are often short and do
not (``where are retries configured?''). Table~\ref{tab:plain} reports five such holdouts. S1, the
short-question fusion of \S\ref{sec:system}, was chosen on the first three sets (looked at), was
neutral on axios, and was confirmed on jsoup, a language and repository never run before:
$+0.333$ R@3 over the packet order and $+0.083$ over BM25. Twelve questions per set is small (one
question is 0.083); we report these as a direction, not a precise effect.

\subsection{Cost}
\label{sec:cost}

% Cost per issue; ours: contributions-log sec. 8.1, sec. 8.2, sec. 8.9, sec. 8.13, sec. 8.18;
% others: related-work-2026.md sec. 1-sec. 2 (SweRank T5, LocAgent T5, Agentless T1).
\begin{table}[t]
\centering
\small
\caption{Cost of localising one issue. Ours measured on a 6-core laptop CPU (Ryzen~5 7530U), cold
index included, no network and no GPU. Others as reported by their authors.}
\label{tab:cost}
\setlength{\tabcolsep}{3pt}
\begin{tabular}{@{}p{0.52\columnwidth}p{0.42\columnwidth}@{}}
\toprule
setting & cost\\
\midrule
NeuroMesh, Lite & p50 2.9\,s, p90 10.6\,s; 0 tokens\\
NeuroMesh, Verified & p50 2.5\,s, p90 8.8\,s; 0 tokens\\
\quad index build, django (3.5k files) & 26\,s cold\\
\quad + dense short-list vote & +$\approx$6\,s p50 (30 definitions + issue)\\
Whole-repository embedding, CPU & 2.6 chunks/s (jina-code-v2)\\
Local 3B LLM judge, CPU & $\approx$105\,s per issue\\
\midrule
SweRankEmbed-L + SweRankLLM-Small (7B) & A100, 12.5\,s, \$0.011\\
LocAgent + Qwen2.5-7B (ft) & GPU, \$0.05\\
LocAgent + Claude-3.5 & \$0.66\\
Agentless + GPT-4o & \$0.70, 78k tokens\\
SWE-Adept + Claude-Sonnet-4.5 & 348k tokens\\
\bottomrule
\end{tabular}
\end{table}


Table~\ref{tab:cost} gives per-issue cost. The engine needs no network, no GPU and no tokens; the
cold index is included in the 2.9\,s median. Its packet (the context an agent receives) has a median
of 11.1k tokens. The cheapest published setting that beats it on file \accat{1} at the same protocol
without an LLM, SweRankEmbed-Small, needs a contrastively trained encoder and an embedding of every
function; at our measured CPU rate for a model of that size, embedding django's roughly 40k chunks
would take hours.

% ---------------------------------------------------------------------------------------------
\section{Negative results and what the holdouts caught}
\label{sec:negative}

% Negative results; contributions-log sec. 8.1, sec. 8.3, sec. 8.4, sec. 8.6, sec. 8.9, sec. 8.13, sec. 8.14, sec. 8.15, sec. 8.16, sec. 8.18.
\begin{table*}[t]
\centering
\small
\caption{Ideas that did not survive the protocol, with the measurement that rejected them. ``dev'' =
SWE-bench dev (tuning split); the bar is file Acc@1 or function Acc@5 $\geq +0.02$ on both halves.}
\label{tab:negative}
\begin{tabular}{@{}p{0.36\textwidth}p{0.20\textwidth}p{0.38\textwidth}@{}}
\toprule
idea (inspiration) & where measured & effect\\
\midrule
Whole-repository dense chunk retrieval on CPU (CodeRankEmbed, Jina) & astroid, 3,364 chunks & 2.6 chunks/s (1,314\,s for one mid-sized repository); hours for django-sized ones\\
Cross-encoder rerank of the top 10 files (jina-reranker-v2) & dev-fast, 59 issues & file Acc@3 0.492 $\to$ 0.373 on long reports; 8.0\,s p50 per issue\\
Same cross-encoder on short questions & ripgrep, click; then fresh cobra & R@3 $+0.375$ / $+0.208$ on the two looked-at sets, $-0.084$ on the fresh one\\
Graph-neighbour vote of the top three files (LocAgent-style expansion without an LLM) & dev-fast & file Acc@1 0.254 $\to$ 0.136\\
Body term frequency in BM25F & dev-fast & file Acc@5 0.544 $\to$ 0.509\\
Commit-history prior (BugLocator, Locus) & dev & file Acc@1 0.271 $\to$ 0.169 (dev-fast) even at weight 0.1\\
Identifiers in the issue's code blocks $\to$ their definitions & dev & file Acc@1 drops; function Acc@5 $+0.019$ (below the bar)\\
RM3 pseudo-relevance feedback on the definition ranking & dev & alone: function Acc@1 0.135 $\to$ 0.038 (dev-fast); as a vote no better than plain BM25\\
Documentation sections as a bridge from question words to identifiers & four plain-language sets & R@3 down on three of four (ripgrep 0.500 $\to$ 0.375, axios 0.833 $\to$ 0.583)\\
Local 3B LLM as file judge (Qwen2.5-Coder-3B, Agentless-style, one call) & 15 dev-fast issues & copies the engine order (0.200/0.467/0.600 unchanged); shuffled input 0.067/0.400/0.600; $\approx$105\,s per issue\\
Learned reranker (logistic regression / boosted trees, 25 features, top 100 functions) & dev halves; dev $\to$ SWE-Gym & function Acc@5 $+0.000$ / $+0.019$ cross-half; $+0.000$ trained on dev, applied to SWE-Gym; the file-rank coefficient flips sign ($-0.61$ vs $+1.23$)\\
\midrule
\multicolumn{3}{@{}l}{\emph{Evaluation slips we made and caught}}\\
Gold match by plain suffix (no path-segment boundary) & cobra & BM25 R@3 printed 0.958 instead of 0.875\\
Function denominator without short reports & Lite 274 & 244 instances and Acc@5 0.443 instead of 274 and 0.394 (v1.3.0)\\
Prototype baseline taken from an older run & dev & every function-level delta re-measured against the true baseline\\
\bottomrule
\end{tabular}
\end{table*}


Table~\ref{tab:negative} lists the ideas that did not survive. Three findings stand out.

\emph{Query length decides whether a cross-encoder helps.} jina-reranker-v2 reordering the engine's
top 10 files improved R@3 on two short-question sets (ripgrep $+0.375$, click $+0.208$) and hurt on
long issue reports (\accat{3} 0.492 $\to$ 0.373): its 512-token window drops the report's specifics,
which the first stage already reads in full. On a fresh short-question set (cobra) it lost $-0.084$;
without the holdout we would have shipped it for short questions.

\emph{Blind graph expansion hurts.} Voting for the graph neighbours of the top three files pulls in
the hubs they call (base classes, utilities) ahead of the edited module (file \accat{1} 0.254
$\to$ 0.136). LocAgent gets value from the graph through an LLM choosing which edges to follow.

\emph{A learned reranker fits the split, not the task.} Logistic regression and boosted trees over 25
features of the top 100 function candidates did not pass cross-half on dev ($+0.000$/$+0.019$);
trained on dev and applied to SWE-Gym they gained nothing, and the largest coefficients change sign
between splits (file rank $-0.61$ on dev, $+1.23$ on SWE-Gym). With a few hundred issues per split,
supervised reranking of lexical features does not generalise.

Two results from the bug-localisation literature were clearly negative on dev. A commit-history prior
in the spirit of BugLocator and Locus (past commits whose message resembles the issue vote for the
files they touched) lowered file \accat{1} from 0.271 to 0.169 on dev-fast even at weight 0.1. RM3
expansion of the definition query was worse than plain BM25 on its own and no better as a vote.
A local 3B LLM (Qwen2.5-Coder-3B-Instruct \citep{hui2024qwen25coder}, 4-bit, llama.cpp
\citep{llamacpp}) as an Agentless-style file judge over the engine's top 10 candidates copied the engine's order when shown it and did worse
when the order was shuffled (\accat{1} 0.200 $\to$ 0.067 on 15 issues), at about 105\,s per issue:
on a laptop the engine's ranking already beats a small LLM judge.

Our own evaluation slips are part of the record: a suffix match without a path-segment boundary that
inflated a baseline, a function-level denominator that silently dropped short reports (244 instead
of 274 instances, \accat{5} 0.443 instead of 0.394), and a prototype baseline taken from an older
run. The denominator slip was caught by re-deriving LocAgent's count of 274; every affected number
was re-measured and is reported corrected.

% ---------------------------------------------------------------------------------------------
\section{Discussion}
\label{sec:discussion}

\paragraph{Where the engine stands.}
It is the best measured point we know of with no learned model at query time, and it is clearly
below trained dense retrievers and LLM pipelines. We see it as a first stage: its top 30 definitions
are a short list that a dense model or an LLM can reorder at a small fraction of the cost of
searching the repository, as the dense vote shows. An LLM stage over \texttt{where\_to\_look},
measured in tokens and \accat{1} against Agentless and LocAgent, is the experiment this paper does
not have (\S\ref{sec:threats}).

\paragraph{Where function misses come from.}
On dev, 45\% of function-level misses have every gold function in the 300-entry list but not in the
top five (median rank 36), 15\% miss the file, and 8\% lack coverage (log~\S8.18). Ordering, not
recall, is the bottleneck --- which is why a reranker of a short list is the natural next stage.

\paragraph{Why holdouts.}
Two of the ideas that looked best on looked-at sets (the cross-encoder on short questions and the
neighbour vote) lost on fresh data, one decided on looked-at sets (S1) was confirmed on a fresh one,
and the gains that did transfer (definition ranking, named files, traceback vote) transferred with
the predicted sign. None of the competitor papers we read reports a pre-registered holdout; on
benchmarks where agent numbers move by 0.09 between re-runs, we think they should.

% ---------------------------------------------------------------------------------------------
\section{Threats to validity}
\label{sec:threats}

\paragraph{Missing LLM stage.}
We do not measure an LLM on top of the engine. A local 3B model was rejected on dev; an API model run
was blocked by unfunded keys at the time of writing. The claim that the engine is a good first stage
for an LLM is supported only by the dense vote and by the 45\% ordering misses.

\paragraph{Tuning-split mismatch.}
Dev and SWE-Gym differ from the holdouts in gold size (3.4--3.7 gold functions vs 1.17 on Lite), so
tuning-split numbers are lower and function-level gains transfer larger than predicted. A split that
resembles Lite would make the bar more sensitive.

\paragraph{Forking paths.}
The dense vote failed the bar on dev and was re-measured on SWE-Gym, a second look at the same idea.
We mitigated this by keeping its weight unchanged, committing the decision before the Lite run, and
reporting the dev result. The effect is modest on each holdout (Lite CI [0.004, 0.051], sign test
$p=0.065$; Verified CI [0.009, 0.048], $p=0.007$), but it has the predicted sign and size on both.

\paragraph{Comparisons across papers.}
All competitor numbers are copied from their papers; our 274 instances are rebuilt from LocAgent's
definition and match its count exactly, but we did not re-run the competitors. Our own BM25 (0.299
file \accat{1}) is below LocAgent's (0.387) on the same instances, so ``gain over BM25'' depends on
which BM25; we report both. Agent numbers vary by model and run.

\paragraph{Small plain-language sets.}
Twelve questions per repository, one annotator, gold written by the authors. MULocBench-style
non-code issues are not covered.

\paragraph{Engineering.}
Timing is from a single Windows laptop shared with other jobs; harness failures (blobless clones,
path-length limits, worktree lock contention) cost instances in early runs and are reported where
they occurred. Benchmark data may appear in the pre-training of the encoders we compare against; the
engine itself has no learned weights, and the dense vote uses an off-the-shelf encoder not trained on
SWE-bench.

% ---------------------------------------------------------------------------------------------
\section{Conclusion}
\label{sec:conclusion}

A model-free engine on a laptop CPU localises SWE-bench issues at file \accat{1} 0.507 (Lite, 276
held-out issues) and 0.435 (Verified outside Lite), more than doubles BM25 on seven repositories it
never saw, and lists every edited function in its top five on 0.460 of LocAgent's Lite instances ---
0.485 when only its top 30 definitions are embedded, on CPU --- in a few seconds and no tokens. It is below trained dense
retrievers and LLM agents. The holdout protocol mattered: two ideas that looked best on looked-at data
lost on fresh data, and most literature priors we tried did not clear a bar fixed before measuring.
What remains: an LLM stage on top of \texttt{where\_to\_look} measured in tokens and accuracy, and
plain-language holdouts of at least fifty questions with a second annotator.

\paragraph{Availability.}
Engine, harness, gold files, result files and the scripts that rebuild every table are open source
\citep{neuromesh}. \todo{archive the release and result files (Zenodo DOI) and cite the exact tag}

\paragraph{Acknowledgements.}
\nm{} is a fork of the open-source NeuroMesh project by Pinoox (\url{https://github.com/pinoox/neuromesh}).
\todo{funding and people}. Development of the engine and of this paper used an LLM coding
assistant; no LLM runs at query time in any reported result. \todo{confirm disclosure wording}

\bibliographystyle{plainnat}
\bibliography{refs}

\end{document}
